\chapter{Projeto 2: Jogo de Perguntas}

\begin{objetivos}
\begin{itemize}
\item Criar um quiz com perguntas e respostas guardadas em listas;
\item Conferir respostas com condições;
\item Mostrar a pontuação final.
\end{itemize}
\end{objetivos}

\section{O plano do quiz}

Neste projeto o gatinho \textbf{faz perguntas} e o jogador digita as respostas. Cada acerto vale um ponto, e no final o gatinho diz quantos pontos você fez.

\begin{definicao}[Roteiro do quiz]
\textbf{Ator:} Gatinho (o apresentador).\\
\textbf{Listas:} \textbf{perguntas} e \textbf{respostas} (a pergunta número 1 combina com a resposta número 1).\\
\textbf{Variáveis:} \textbf{pontos} e \textbf{número}.
\end{definicao}

\begin{missao}[Missão 18: Monte o Quiz]
Crie um quiz com 3 perguntas sobre o que quiser, usando listas.
\end{missao}

\section{Passo 1: preparando as listas}

Crie as listas \textbf{perguntas} e \textbf{respostas}. No começo do jogo, apague tudo e cadastre as perguntas. A ordem importa!

\begin{codigo}[Cadastrando as perguntas]
\begin{pilha}
\bk[scEvents]{quando \bandeira\ for clicado}
\bk[scLists]{apague tudo de \menu{perguntas}}
\bk[scLists]{apague tudo de \menu{respostas}}
\bk[scLists]{adicione \entrada{Quanto é 2 + 2?} a \menu{perguntas}}
\bk[scLists]{adicione \entrada{4} a \menu{respostas}}
\bk[scLists]{adicione \entrada{Qual é a capital do Brasil?} a \menu{perguntas}}
\bk[scLists]{adicione \entrada{Brasília} a \menu{respostas}}
\bk[scLists]{adicione \entrada{Quantos dias tem a semana?} a \menu{perguntas}}
\bk[scLists]{adicione \entrada{7} a \menu{respostas}}
\end{pilha}
\end{codigo}

\section{Passo 2: fazendo as perguntas}

Para cada pergunta, o gatinho pergunta, espera a resposta e compara com a resposta certa, que está \textbf{na mesma posição} da outra lista.

Encaixe os blocos abaixo na mesma pilha da bandeira, logo depois de cadastrar a última resposta. Assim, o quiz só começa depois que as duas listas estiverem prontas.

\begin{codigo}[O jogo]
\renewcommand{\sbw}{410pt}
\begin{pilha}
\bk[scVariables]{defina \menu{pontos} para \entrada{0}}
\bk[scVariables]{defina \menu{número} para \entrada{1}}
\sbc[scControl]{repita \reporter[scLists]{tamanho de \menu{perguntas}} vezes}
\bk[scSensing]{pergunte \reporter[scLists]{item \variavel{número} de \menu{perguntas}} e espere}
\sbc[scControl]{se \booleano{\variavel{resposta} = \reporter[scLists]{item \variavel{número} de \menu{respostas}}} então}
\bk[scLooks]{diga \entrada{Acertou!} por \entrada{1} segundos}
\bk[scVariables]{mude \menu{pontos} em \entrada{1}}
\sbcelse{senão}
\bk[scLooks]{diga \entrada{Errou!} por \entrada{1} segundos}
\sbcend
\bk[scVariables]{mude \menu{número} em \entrada{1}}
\sbcend
\bk[scLooks]{diga \reporter{junte \entrada{Você fez } \reporter{junte \variavel{pontos} \entrada{ pontos!}}} por \entrada{4} segundos}
\end{pilha}
\end{codigo}

\begin{dica}
Esse código funciona com \textbf{qualquer quantidade} de perguntas! O laço repete o número de vezes que a lista tem. Para ter mais perguntas, basta cadastrar mais.
\end{dica}

\begin{cuidado}
Se o jogador digitar ``brasília'' em minúsculas, o Scratch ainda considera igual a ``Brasília'' (ele não diferencia maiúsculas). Mas acentos e espaços a mais contam como erro. Combine isso com os jogadores!
\end{cuidado}

\section{Melhorando o quiz}

\begin{exercicio}[Mais uma pergunta]
Adicione uma quarta pergunta ao quiz. Do que você precisa mudar?
\end{exercicio}

\begin{solucao}
Acrescente a nova pergunta e a resposta ao fim do cadastro, antes do trecho O jogo:
\begin{pilha}
\bk[scLists]{adicione \entrada{Quantos lados tem um triângulo?} a \menu{perguntas}}
\bk[scLists]{adicione \entrada{3} a \menu{respostas}}
\end{pilha}
O laço já repete o tamanho da lista, então o resto do código não muda.
\end{solucao}

\begin{exercicio}[Mensagem de vitória]
No final, se o jogador acertar \textbf{todas}, o gatinho diz ``Parabéns, você é demais!''. Senão, diz ``Tente de novo!''.
\end{exercicio}

\begin{solucao}
\begin{pilha}
\sbc[scControl]{se \booleano{\variavel{pontos} = \reporter[scLists]{tamanho de \menu{perguntas}}} então}
\bk[scLooks]{diga \entrada{Parabéns, você é demais!}}
\sbcelse{senão}
\bk[scLooks]{diga \entrada{Tente de novo!}}
\sbcend
\end{pilha}
\end{solucao}

\begin{exercicio}[Som de acerto e erro]
Toque um som alegre quando o jogador acertar e um som triste quando errar.
\end{exercicio}

\begin{solucao}
Coloque um bloco \bloco[scSound]{inicie o som} dentro de cada ramo do \bloco[scControl]{se}, junto com o \bloco[scLooks]{diga}.
\end{solucao}

\begin{desafio}[Placar visual]
Faça um segundo ator que muda de fantasia conforme os pontos: triste com 0, alegre a partir de 3. Dica: use a mensagem ``atualizou'' sempre que os pontos mudarem.
\end{desafio}

\begin{solucao}
No quiz, envie \bloco[scEvents]{envie \menu{atualizou}} em dois lugares: logo depois de zerar \textbf{pontos} e logo depois de somar 1. No segundo ator:
\begin{pilha}
\bk[scEvents]{quando eu receber \menu{atualizou}}
\sbc[scControl]{se \booleano{\variavel{pontos} $>$ \entrada{2}} então}
\bk[scLooks]{mude para a fantasia \menu{alegre}}
\sbcelse{senão}
\bk[scLooks]{mude para a fantasia \menu{triste}}
\sbcend
\end{pilha}
\end{solucao}
